\documentclass[10pt,a4paper]{amsart}
\usepackage{amsmath,amssymb,amsthm,graphicx}
\usepackage{verbatim}
\newenvironment{reference}{\comment}{\endcomment}
%\newenvironment{reference}{}{}
\newenvironment{slogan}{\comment}{\endcomment}
\newenvironment{history}{\comment}{\endcomment}

% For commutative diagrams we use Xy-pic
\usepackage[all]{xy}

% We use 2cell for 2-commutative diagrams.
\xyoption{2cell}
% Theorem environments.
%
\theoremstyle{plain}
\newtheorem{theorem}[subsection]{Theorem}
\newtheorem{proposition}[subsection]{Proposition}
\newtheorem{lemma}[subsection]{Lemma}

\theoremstyle{definition}
\newtheorem{definition}[subsection]{Definition}
\newtheorem{example}[subsection]{Example}
\newtheorem{exercise}[subsection]{Exercise}
\newtheorem{situation}[subsection]{Situation}

\theoremstyle{remark}
\newtheorem{remark}[subsection]{Remark}
\newtheorem{remarks}[subsection]{Remarks}

\numberwithin{equation}{subsection}

% Macros
%
\def\lim{\mathop{\mathrm{lim}}\nolimits}
\def\colim{\mathop{\mathrm{colim}}\nolimits}
\def\Spec{\mathop{\mathrm{Spec}}}
\def\Hom{\mathop{\mathrm{Hom}}\nolimits}
\def\Ext{\mathop{\mathrm{Ext}}\nolimits}
\def\SheafHom{\mathop{\mathcal{H}\!\mathit{om}}\nolimits}
\def\SheafExt{\mathop{\mathcal{E}\!\mathit{xt}}\nolimits}
\def\Sch{\mathit{Sch}}
\def\Mor{\mathop{\mathrm{Mor}}\nolimits}
\def\Ob{\mathop{\mathrm{Ob}}\nolimits}
\def\Sh{\mathop{\mathit{Sh}}\nolimits}
\def\NL{\mathop{N\!L}\nolimits}
\def\CH{\mathop{\mathrm{CH}}\nolimits}
\def\proetale{{pro\text{-}\acute{e}tale}}
\def\etale{{\acute{e}tale}}
\def\QCoh{\mathit{QCoh}}
\def\Ker{\mathop{\mathrm{Ker}}}
\def\Im{\mathop{\mathrm{Im}}}
\def\Coker{\mathop{\mathrm{Coker}}}
\def\Coim{\mathop{\mathrm{Coim}}}

% Boxtimes
%
\DeclareMathSymbol{\boxtimes}{\mathbin}{AMSa}{"02}

%
% Macros for moduli stacks/spaces
%
\def\QCohstack{\mathcal{QC}\!\mathit{oh}}
\def\Cohstack{\mathcal{C}\!\mathit{oh}}
\def\Spacesstack{\mathcal{S}\!\mathit{paces}}
\def\Quotfunctor{\mathrm{Quot}}
\def\Hilbfunctor{\mathrm{Hilb}}
\def\Curvesstack{\mathcal{C}\!\mathit{urves}}
\def\Polarizedstack{\mathcal{P}\!\mathit{olarized}}
\def\Complexesstack{\mathcal{C}\!\mathit{omplexes}}
% \Pic is the operator that assigns to X its picard group, usage \Pic(X)
% \Picardstack_{X/B} denotes the Picard stack of X over B
% \Picardfunctor_{X/B} denotes the Picard functor of X over B
\def\Pic{\mathop{\mathrm{Pic}}\nolimits}
\def\Picardstack{\mathcal{P}\!\mathit{ic}}
\def\Picardfunctor{\mathrm{Pic}}
\def\Deformationcategory{\mathcal{D}\!\mathit{ef}}
\usepackage{cleveref}
\crefname{originallemma}{original lemma}{original lemmas}
\Crefname{originallemma}{Original lemma}{Original lemmas}
\makeatletter
\newlabel{coherent-lemma-flat-base-change-cohomology}{{02KH}{source}{}{sourcepointer-1}{}}
\newlabel{coherent-lemma-flat-base-change-cohomology@cref}{{[originallemma][0][]02KH}{[1][0][]source}{}{}{}}
\makeatother
\crefname{originallemma}{original lemma}{original lemmas}
\Crefname{originallemma}{Original lemma}{Original lemmas}
\makeatletter
\newlabel{divisors-lemma-Pic-projective-space-UFD}{{0BXJ}{source}{}{sourcepointer-2}{}}
\newlabel{divisors-lemma-Pic-projective-space-UFD@cref}{{[originallemma][0][]0BXJ}{[1][0][]source}{}{}{}}
\makeatother
\crefname{originalreference}{original reference}{original references}
\Crefname{originalreference}{Original reference}{Original references}
\makeatletter
\newlabel{lemma-blowup-pic}{{0C5G}{source}{}{sourcepointer-3}{}}
\newlabel{lemma-blowup-pic@cref}{{[originalreference][0][]0C5G}{[1][0][]source}{}{}{}}
\makeatother
\crefname{originalreference}{original reference}{original references}
\Crefname{originalreference}{Original reference}{Original references}
\makeatletter
\newlabel{lemma-cohomology-of-blowup}{{0AGS}{source}{}{sourcepointer-4}{}}
\newlabel{lemma-cohomology-of-blowup@cref}{{[originalreference][0][]0AGS}{[1][0][]source}{}{}{}}
\makeatother
\title{The Picard group of the contraction of an exceptional curve of the first kind}
\author{The Stacks Project Authors}
\date{stacks-project:resolve.tex:local-snapshot-20261012; GNU FDL 1.2 or later}
\begin{document}
\maketitle
\noindent\textit{Editorial note.} All delivered mathematical text is the author's own, copied from the pinned Stacks Project chapter file resolve.tex; no statement or proof step was rewritten, reordered or removed.\par\medskip
\noindent\textit{Editorial note.} The counted proof is the author's complete argument: the identification of the restriction map with the degree and the reduction to the statement on Picard groups, the verification that the pullback is an isomorphism onto the invertible sheaves restricting trivially to the curve, and the verification that the pushforward inverts the pullback after reduction to a regular local base, together with the explicit splitting.\par\medskip
\noindent\textit{Editorial note.} The label lemma-pic-blowup names the counted statement, and the proof environment immediately following it in the source is delivered verbatim as its proof.\par\medskip
\noindent\textit{Editorial note.} The counted statement is carried by the author's own lemma environment in the chapter source.\par\medskip
\noindent\textit{Editorial note.} Pre-open gate: python operations/agent-authoring/tools/precard\_gate.py --source resolve.tex --label lemma-pic-blowup --start 4776 --end 4836 --local-id 28086 returned RESULT=OK (L28086 free, occupancy FREE, public interval clear). Pre-assembly tag self-check: canonical tag resolved and not present in reference/public-items.jsonl.\par\medskip
\noindent\textit{Editorial note.} The statement and proof are copied from the pinned chapter file; the cited inputs are kept as external pointers with their canonical Stacks tags.\par\medskip
\noindent\textit{Allowed standard prerequisites (editorial).}\begin{itemize}
\item The Picard group of the projective line and the degree of the exceptional divisor. Divisors, the lemma computing the Picard group of a projective space, together with the defining property of an exceptional curve of the first kind that the restriction of its associated invertible sheaf has degree minus one. Supplies the identification of the last map with restriction to the curve and the computation of the degree of the restriction of the associated invertible sheaf.
\item Cohomology of the structure sheaf of the blowup and triviality of invertible sheaves over a regular local base. Same chapter, the lemma computing the cohomology of the structure sheaf of the blowup and, in its last part, the statement that the pushforward of the structure sheaf is the structure sheaf, together with the lemma describing the Picard group of the blowup of a regular local ring. Supplies the isomorphism between a module on the contraction and the pushforward of its pullback, and the triviality of invertible sheaves after restricting to the local ring of the contraction at the image point.
\item Flat base change for pushforward, pullback and the adjunction map. Cohomology of Schemes, the flat base change lemma, applied to the computation of the pushforward and of the adjunction map. Supplies the reduction of the two claims to the local ring of the contraction at the image point of the curve.
\end{itemize}\medskip
\ifdefined\hypertarget\hypertarget{sourcepointer-1}{}\fi
\noindent\textit{Source reference.} Original lemma 02KH: sources/stacks-chow/tags.txt, source lines 3372--3372. Prerequisite pointer invoked at proof line(s) 4821; it names a standard fact of a sibling chapter that the counted argument uses, and every step of the counted proof built on it is retained verbatim in the delivered body.\par\medskip
\ifdefined\hypertarget\hypertarget{sourcepointer-2}{}\fi
\noindent\textit{Source reference.} Original lemma 0BXJ: sources/stacks-chow/tags.txt, source lines 14982--14982. Prerequisite pointer invoked at proof line(s) 4793; it names a standard fact of a sibling chapter that the counted argument uses, and every step of the counted proof built on it is retained verbatim in the delivered body.\par\medskip
\ifdefined\hypertarget\hypertarget{sourcepointer-3}{}\fi
\noindent\textit{Source reference.} Original reference 0C5G: sources/stacks-resolve/pinned-original-source.tex, source lines 756--756. Prerequisite pointer invoked at proof line(s) 4825; it names a standard fact of the same chapter that the counted argument uses, and every step of the counted proof built on it is retained verbatim in the delivered body.\par\medskip
\ifdefined\hypertarget\hypertarget{sourcepointer-4}{}\fi
\noindent\textit{Source reference.} Original reference 0AGS: sources/stacks-resolve/pinned-original-source.tex, source lines 786--786. Prerequisite pointer invoked at proof line(s) 4809; it names a standard fact of the same chapter that the counted argument uses, and every step of the counted proof built on it is retained verbatim in the delivered body.\par\medskip
\begin{lemma}
\label{lemma-pic-blowup}
Let $b : X \to X'$ be the contraction of an
exceptional curve of the first kind $E \subset X$.
Then there is a short exact sequence
$$
0 \to \Pic(X') \to \Pic(X) \to \mathbf{Z} \to 0
$$
where the first map is pullback by $b$ and the second map sends
$\mathcal{L}$ to the degree of $\mathcal{L}$ on the exceptional
curve $E$. The sequence is split by the map
$n \mapsto \mathcal{O}_X(-nE)$.
\end{lemma}

\begin{proof}
Since $E = \mathbf{P}^1_k$ we see that the Picard group of $E$
is $\mathbf{Z}$, see Divisors, Lemma
\ref{divisors-lemma-Pic-projective-space-UFD}.
Hence we can think of the last map as $\mathcal{L} \mapsto \mathcal{L}|_E$.
The degree of the restriction of $\mathcal{O}_X(E)$ to $E$ is $-1$
by definition of exceptional curves of the first kind. Combining these
remarks we see that it
suffices to show that $\Pic(X') \to \Pic(X)$ is injective
with image the invertible sheaves restricting to $\mathcal{O}_E$ on $E$.

\medskip\noindent
Given an invertible $\mathcal{O}_{X'}$-module
$\mathcal{L}'$ we claim the map $\mathcal{L}' \to b_*b^*\mathcal{L}'$
is an isomorphism. This is clear everywhere except possibly at the image
point $x \in X'$ of $E$. To check it is an isomorphism on stalks
at $x$ we may replace $X'$ by an open neighbourhood at $x$ and
assume $\mathcal{L}'$ is $\mathcal{O}_{X'}$. Then we have to
show that the map $\mathcal{O}_{X'} \to b_*\mathcal{O}_X$
is an isomorphism. This follows from Lemma \ref{lemma-cohomology-of-blowup}
part (4).

\medskip\noindent
Let $\mathcal{L}$ be an invertible $\mathcal{O}_X$-module with
$\mathcal{L}|_E = \mathcal{O}_E$. Then we claim 
(1) $b_*\mathcal{L}$ is invertible and
(2) $b^*b_*\mathcal{L} \to \mathcal{L}$ is an isomorphism.
Statements (1) and (2) are clear over $X' \setminus \{x\}$.
Thus it suffices to prove (1) and (2) after base change
to $\Spec(\mathcal{O}_{X', x})$.
Computing $b_*$ commutes with flat base change
(Cohomology of Schemes, Lemma \ref{coherent-lemma-flat-base-change-cohomology})
and similarly for $b^*$ and formation of the adjunction map.
But if $X'$ is the spectrum of a regular local ring
then $\mathcal{L}$ is trivial by the description of
the Picard group in Lemma \ref{lemma-blowup-pic}. Thus
the claim is proved.

\medskip\noindent
Combining the claims proved in the previous two paragraphs we
see that the map $\mathcal{L} \mapsto b_*\mathcal{L}$
is an inverse to the map
$$
\Pic(X') \longrightarrow \Ker(\Pic(X) \to \Pic(E))
$$
and the lemma is proved.
\end{proof}

\end{document}
